\section{Implementation and Evaluation}
\subsection{Implementation Details}
  %1. system components
\noindent{\bf System Implementation}
  The \ttt{ZkregPlus} system contains three parts: (1)
a parser of ClamAV regex signatures, (2) the FoldPot 
framework which provides folding and zkSNARK, and
(3) a collection of circuits that implements 
streamed CP, SED, and DFA discharging algorithms.
The system is implemented in Rust, consisting of 
around {\em to check (30k)} lines of code, and its source
code is publicaly available on GitHub.
\footnote{\url{http://github.com/xxxx}}

 The regex parser supports ClamAV's own hex signature format,
and as well as full PCRE regex specificaton. We rely on the
\ttt{regex\_syntax} crate of Rust to parse PCRE regex into
a high-level intermediate representation, from where, either
AC-DFAs are constructed for sub-signatures, or 
DFA built using the \ttt{Rustomaton} crate\footnote{\url{https://crates.io/crates/rustomaton}}, or SED related information are generated.
The signature related information (e.g., transition table)
are encoded into a two-column global lookup table.
Table\ \ref{tbl:lkup} shows
the break-down of the lookup table which has
$2^{28}$ entries.

\begin{small}
\begin{table}
\begin{tabular}{|l|c|}
  \hline
  Sub-Lookup Table &  Size \\
  \hline
  (1) range table1 $[0,16]$ & $16$ \\ 
  (2) range table2 $[0,2^{23}]$ & $16$ \\ 
  (3) AC-DFA for CP (regular) & $2000$ \\ 
  (4) AC-DFA for CP (ignore case) & $2000$ \\
  (5) AC-DFA for SED (regular) & $2000$ \\ 
  (6) AC-DFA for SED (ignore case) & $2000$ \\
  (7) 16 (hard) DFAs  & $2000$ \\
  (8) signature DB (e.g., subsig-step-pattern
  	relations) & $2000$ \\
  (9) tri-logic computation table & $2000$ \\
  \hline
\end{tabular}
\caption{Breakdown of 2-column Lookup Table (TO FILL DETAILS)}
\label{tbl:lkup}
\end{table}
\end{small}

The FoldPot framework is built over the
Sonobe framework \cite{Sonobe}, extending its Nova + Cyclefold framework
to SuperNova \cite{SuperNova} and then we implement other
components such as QA-NIZK \cite{Gonzalez15},
the Sigma-IR1CS protocol, scheme for distributing
Logup lookup \cite{Hab22} checks into shares as suggested
in Mangrove \cite{Mangrove}, and the 
cycle pairing framework as shown in Section \ref{sec:cyclepair}.
We use Bn254 and Grumpkin curves \cite{Grumpkin} 
(s.t., Grumpkin's scalar field is the base field of Bn254),
and instantiate hash functions using Poseidon \cite{Poseidon}.
Most crypto primitives and polynomial related functions
are built over the Arkworks library \cite{arkworks}.
We further developed an additional layer of
relational database operations to support
CP, SED, and DFA discharging algorithms,
which heavily rely on lookups.

\noindent{\bf Experimental Setup} All of our experiments are run on
a Google Cloud Computing (GCP) xxyyzz model server with xxyyzz cores
and XXYYZZ GB of RAM. The folding stages needs XXXXYY GB, and
decider consumes all XXYYZZGB of RAM.

\noindent{\bf Evaluation Questions} There are two basic essential
research questions to answer: 
\begin{itemize}
	\item $\mathbf{Q_1}$: How effective is the approximation
	technique? In particular, What is the false positives generated
	by approximations? 

	\item $\mathbf{Q_2}$: What is the performance of the 
	proposed approach in terms of prover speed and proof size?
\end{itemize}

\subsection{Evaluation of Approximation}
To answer $\mathbf{Q_1}$, we evaluate the approximation algorithms
using the real data set of entire ClamAV logical signatures.
The set is comprised of $38876$ signatures.

Fourteen signatures (14)  are dropped
from the original $38889$ signatures,
because after applying approximation
they report false positives on
some Linux executables or Enron emails which are benign. 
By ``benign" we mean: when running clamav using the specific rule set,
ClamAV does not report positive on matching virus signatures.
The breakdown of these false positives are listed below.
(i) Over approximation of enlarged target field.
We do not handle the \texttt{Target} meta-data, which specifies
to which categories of executables that the signature should apply to.
In our approximation, we just ignore the target field and apply
it all all.  \texttt{Win.Trojan.FakeSkype} was targeted at Win32 PE,
and it raises false positive \texttt{vmlinuz} (Linux kernel image). 
(ii) Over approximation  program entry point position
constraint ignored.
For example, in \texttt{Win.Trojan.Generic-43;}
\texttt{...;EP+0:5250{-8}5133c0;...},
the position constraint needs to locate executable file entry point.
This exceeds regex ability, we approximate it
with \texttt{EP+0:} with wildcard string \texttt{.*} before the patterns. 
However, such patterns are too short and
causing false positives. 
(iii) Over approximation of multiline mode.  We use multiline mode 
for all signatures, i.e., new line character is
matched by wildcard dot. In this case, one signature (which works in
single line mode) is removed
for causing false positive due to multi-line approximation.
(iv) Both DFA and SED are expensive. This applies to 
\ttt{Js.File.MaliciousHeuristic-6260249-0}
where its signature is a concatenation of patterns 
"\ttt{...s(?P=junk)}\texttt{c(?P=}\texttt{junk)}\\ \texttt{r(?P=junk)}\texttt{i(?P=junk)...}", 
where for SED, the pattern words are too short and it is too expensive
to construct its DFA.

Among all remaining $38875$ ($99.96\%$) of the $38889$ signatures,
the vast majority have hex subsignatures and $559$ of them have
PCRE subsignatures. Six ($6$) of these PCRE signatures use features such
as named back-references and look-arounds, and they are approximated.

We use two sets of target data set (to discharge free of malware). The first 
is a set of all binary executable files in a Linux CentOS ($2702$ files
total $752$ MB) and Enron Email Data Set ($517401$ files totally $2.74$ GB).
Running ClamAv to discharge all binary executables takes about
$10$ seconds, but much longer ($4$ hours) on Enron Data 
set (mainly due to file reading cost).

\begin{figure}
\begin{footnotesize}
\begin{center}
\begin{tabular}{c c c c c c c}
\hline
\multicolumn{7}{c}{\bf ClamAV Logical Signatures (38886) TO UPDATE!!!} \\
\hline
\multicolumn{7}{c}{Centos Binary Executables (2702 samples)} \\
\hline
log(file) & files & After $\cp \Rightarrow$ &  $\sde \Rightarrow$ & 
	$\dfa \Rightarrow$  \\
\hline
\hline
7 &   72     & (87,90,72)  & (0,0,0) & (0,0,0) \\
8 &   0  & (0,0,0) & (0,0,0) & (0,0,0) \\
9 &   0  & (0,0,0) & (0,0,0) & (0,0,0) \\
10&   2  & (96,101,2)  & (0,0,0) & (0,0,0) \\
11&   13     & (96,103,13) & (0,0,0) & (0,0,0) \\
12&   16     & (98,105,16) & (0,0,0) & (0,0,0) \\
13&   105    & (98,105,105)    & (0,0,0) & (0,0,0) \\
14&   727    & (105,119,727)   & (0,1,2) & (0,0,0) \\
15&   478    & (104,162,478)   & (0,2,1) & (0,0,0) \\
16&   442    & (107,163,442)   & (0,2,70)    & (0,0,0) \\
17&   273    & (116,164,273)   & (0,2,6) & (0,0,0) \\
18&   163    & (122,183,163)   & (0,2,7) & (0,0,0) \\
19&   158    & (131,191,158)   & (0,2,60)    & (0,0,0) \\
20&   67     & (145,201,67)    & (1,3,55)    & (0,0,0) \\
21&   153    & (185,202,153)   & (1,3,146)   & (0,0,0) \\
22&   14     & (150,173,14)    & (1,3,11)    & (0,0,0) \\
23&   6  & (175,235,6) & (3,4,6) & (0,0,0) \\
24&   9  & (224,277,9) & (2,4,9) & (0,1,1) \\
25 &  3  & (226,270,3) & (3,4,3) & (0,1,2) \\
26 &  0  & (0,0,0) & (0,0,0) & (0,0,0) \\
27 &  1  & (285,285,1) & (4,4,1) & (0,0,0) \\
\multicolumn{7}{c}{Enron Email Data (517401 samples)} \\
\hline
9 &   14025  & (87,147,14025)  & (0,0,0) & (0,0,0) \\
10 &   141762     & (88,148,141762) & (0,0,0) & (0,0,0) \\
11 &  175664     & (90,156,175664) & (0,0,0) & (0,0,0) \\
12 &  113795     & (91,160,113795) & (0,0,0) & (0,0,0) \\
13 &  47882  & (93,166,47882)  & (0,0,0) & (0,0,0) \\
14 &  17421  & (96,172,17421)  & (0,0,0) & (0,0,0) \\
15 &  4681   & (98,162,4681)   & (0,0,0) & (0,0,0) \\
16 &  1331   & (101,154,1331)  & (0,0,0) & (0,0,0) \\
17 &  644    & (104,164,644)   & (0,0,0) & (0,0,0) \\
18 &  175    & (110,177,175)   & (0,0,0) & (0,0,0) \\
19 &  12     & (113,119,12)    & (0,0,0) & (0,0,0) \\
20 &  3  & (107,108,3) & (0,0,0) & (0,0,0) \\
21 &  6  & (105,110,6) & (0,0,0) & (0,0,0) \\
\end{tabular}
\end{center}
\end{footnotesize}
\caption{Preliminary Data}
\footnotesize{Column 1: the log of file size. Column 2: the number of file samples.
For the rest of columns, each tuple (avg, max, count) represents the average and max
number of signatures that can not be discharged by the approach. The count represents
the number of files that need to go to next discharging stage.}
\label{fig:predata}
\end{figure}

Given the cost of the proposed schemes, we discharge an input file following
the order of $\cp$, $\bw$, $\sde$, $\dfa$, and $\isde$. Once a signature
is discharged, it does not go to the next stage. At the last stage
a count $0$ means that a file is completely discharged of all virus signatures
(thus malware free).  {\bf CHECK:} We summarize the asymptotic complexity of
these methods.  Let $u$ be the number of signatures to discharge
at a stage, their prover complexity are:
$O(|s|)$, $O(|s|)$,  $O(|s|)$, $O(u \times |s|)$, $O(u \times |s|)$, respectively.
Note that for the first $3$ stages, the prover complexity is
{\em independent} of the number of signatures to discharge.
In our real data set, the $u$ for the last two stages is less than $2$.


In Figure\ \ref{fig:predata} we display the effectiveness of the proposed
$3$-stage discharging scheme, over the binary executable data set and
Enron data set. The Enron set performs similarly, while it relies more
on $\sde$ because Enron dataset are mostly alpha-numerical and there are 
many more ``false positives" on patterns if no checks on structures are
performed.

Take the $2^{20}$ file size category in CentOS dataset as an example:
\[
\begin{footnotesize}
\begin{tabular}{c c c c c c c}
log(file) & files & After $\cp \Rightarrow$ & $\sde \Rightarrow$ & 
	$\dfa \Rightarrow$ \\
\hline
20  &   67       &  (143,198,67)       &  (1,2,54)      &  (0,0,0)   \\
\end{tabular}
\end{footnotesize}
\]

There are $67$ files in the category of $1MB$ size. The first $\cp$ stages
screens the vast majority ($99.5\%$) of signatures and in average there are
$143$ (max $198$) out of $38k$ signatures left. The next $\bw$ stages
discharged in average $39$ signatures, and the $\sde$ discharged all of
the rest except $1$ (in average) and $2$ (max) for $54$ files. 
None of the signatures for $1MB$ files would have to go through to the
last $\isde$ stage. For the entire CentOS dataset,
in the worst case, we only have
to discharge $7$ signatures via $\dfa$ for one file, however, this happens
very rarely. For both the CentOS and Enron dataset,
there are only two signatures\footnote{$\texttt{Js.File.MaliciousHeurstic-6260249-0}$ and $\texttt{Html.Exploit.CVE\_2017\_0069}$ $\texttt{-6013012-3}$}
that have to to through $\isde$ for $5$ (out of $2702$) binary executable files
and $6$ (out of $517401$) Enron email files.
In summary, given that for the last two stages the number of signatures is less than
$2$, the total prover cost is just a constant factor ($\leq 5$) of the
input string size, i.e.,
\[
	5 \times |s|.
\]
This is a speed-up of $7.7\times 10^{6}$
compared with the state of the art (\cite{Zombie23,Luo23})\footnote{In the case of Reef \cite{Reef} it works very well for the non-membership proof case 
where the patterns start
at fixed positions. However, if there is a \texttt{.*} preceding regex patterns,
our technique works better. It is not known if Reef's alternating automata 
can be extended easily to mass process large collection of signatures.}, which
use NFA discharging each
signature separately, where the average $|\texttt{NFA}|>1000$
(given average length of ClamAV signature is $300$):
\[
38878 \times |s| \times |\texttt{NFA}|.
\]

The speed-up comes from two sources. First, we ``mass-process" large collections
of signatures using ACDFA in the first three stages. Second, we rely on
pre-processed lookup arguments (whose prover complexity is independent
of the lookup/container table size) in proving DFA acceptance, thus eliminating
the size of automata in the prover cost. Several different types of
look-up arguments in used in our zk-prover scheme, which is demonstrated
in details in Section 4.

\subsection{Overall System Performance}
We then try to answer the question what is the system performance?

(1) micro-bench mark: group and field operations (how many CPUs).
(2) data size: all 2700 binary samples, chunked into 128k. 
  i) Show the variety of configs for circuit setup (like different
  	size for how many dfas, and the size of step-pattern stack).
  ii) Show an example of a sample circuit cost breakdown. cost of
  	lkups, cost of circuit itself, and cost of support SED.
(3) folding speed
(4) decider step cost. 
	i) show the breakdown of the decider circuit.
	ii) report the time and space consumption. 
	iii) report the proof size breakdown.


\subsection{Micro Benchmarks}

\subsection{Learn from Other Papers}
(1) ZkMiddleBox: eval questions: baseline perf and optimization; case study,
cost for real-world case. http firewalling, domain blocklisting.
no comparative work.

(2) Zombie: eval questions: (1) overheads added by config of Zombie; (2) 
how does batching improve throughput? (3) costs of different components?
(4) how effective are zombie regex?A

(3) Reef: overal performance (compilation, solving, proving, ver)A
 comparaive performance, different appraoches (dfa, dfa + recursion, safa
 	+ lookup), 

